\chapter{Perancangan}
\label{chap:perancangan}

\section{Perancangan Perangkat Lunak}
\label{sec:perancangan-pl}

\subsection{Desain Kelas}
\label{subsec:desain-kelas}

Sistem dibangun menggunakan arsitektur modular yang memisahkan tiga tanggung jawab utama: pengelolaan data papan, logika \textit{solver}, dan antarmuka pengguna. Pemisahan ini memungkinkan kedua \textit{solver} beroperasi secara independen terhadap representasi visual dan sebaliknya, sehingga perubahan pada salah satu bagian tidak berdampak pada bagian lainnya. Diagram kelas yang menunjukkan hubungan antarkeseluruhan kelas dalam sistem diilustrasikan pada Gambar~\ref{fig:arsitektur}.

\begin{figure}[!htbp]
	\centering
	\includegraphics[width=0.62\linewidth]{cDiagram.png}
	\caption{Diagram kelas sistem yang menunjukkan hubungan antara kelas data, kelas \textit{solver}, dan kelas antarmuka.}
	\label{fig:arsitektur}
\end{figure}

Penjelasan rinci untuk setiap kelas, mencakup kegunaan, atribut, dan metode, dijabarkan pada bagain \ref{subsubsec:kelas-cell}. Selain kelas-kelas tersebut, sistem juga memiliki dua titik masuk (\textit{entry point}) yang terpisah. Titik masuk pertama adalah kelas \texttt{Main} yang menginisialisasi antarmuka pengguna dengan membuat \textit{instance} dari \texttt{Homepage}, sehingga pengguna dapat memainkan \textit{Colored Queens} melalui antarmuka grafis. Titik masuk kedua adalah kelas \texttt{BatchMain} yang digunakan secara khusus untuk keperluan eksperimen: kelas ini menjalankan kedua \textit{solver} pada seluruh dataset puzzle secara otomatis dan paralel, kemudian mencatat metrik kinerja ke dalam berkas \texttt{.txt}. Kedua kelas ini tidak dijabarkan secara rinci karena hanya berfungsi sebagai pengelola yang memanggil kelas-kelas lain tanpa memiliki atribut atau relasi yang perlu dimodelkan. Rincian implementasi \texttt{BatchMain} dibahas pada Bagian~\ref{subsec:impl-pengujian}.

\subsubsection{Kelas \texttt{Cell}}
\label{subsubsec:kelas-cell}

Kelas \texttt{Cell} merepresentasikan satu sel pada papan permainan \textit{Colored Queens}. Setiap objek \texttt{Cell} menyimpan posisi sel dalam koordinat baris dan kolom serta informasi warna dalam format RGB, sehingga kelas ini berfungsi sebagai satuan data terkecil yang digunakan untuk membangun struktur papan.

Atribut-atribut kelas \texttt{Cell} dirangkum pada Tabel~\ref{tab:attr-cell}.

\begin{longtable}{llp{7.5cm}}
	\caption{Atribut kelas \texttt{Cell}}
	\label{tab:attr-cell} \\
	\toprule
	\textbf{Atribut} & \textbf{Tipe} & \textbf{Deskripsi} \\
	\midrule
	\endfirsthead
	\toprule
	\textbf{Atribut} & \textbf{Tipe} & \textbf{Deskripsi} \\
	\midrule
	\endhead
	\bottomrule
	\endfoot
	\texttt{row} & \texttt{int} & Koordinat baris sel pada papan (indeks berbasis nol) \\
	\midrule
	\texttt{col} & \texttt{int} & Koordinat kolom sel pada papan (indeks berbasis nol) \\
	\midrule
	\texttt{r} & \texttt{int} & Komponen \textit{red} dari warna sel (0--255) \\
	\midrule
	\texttt{g} & \texttt{int} & Komponen \textit{green} dari warna sel (0--255) \\
	\midrule
	\texttt{b} & \texttt{int} & Komponen \textit{blue} dari warna sel (0--255) \\
	\bottomrule
\end{longtable}

Kelas ini menyediakan \textit{constructor} yang menerima kelima atribut di atas, metode \textit{getter} dan \textit{setter} untuk setiap atribut, serta metode \texttt{toString()} yang mengembalikan representasi tekstual sel dalam format \texttt{Cell[row=\_, col=\_, rgb=(\_,\_,\_)]} untuk keperluan \textit{debugging}.

\subsubsection{Kelas \texttt{Board}}
\label{subsubsec:kelas-board}

Kelas \texttt{Board} merepresentasikan keseluruhan papan permainan \textit{Colored Queens}. Kelas ini bertanggung jawab mengagregasi seluruh objek \texttt{Cell} ke dalam sebuah pemetaan yang mengelompokkan sel-sel berdasarkan warna, sehingga \textit{solver} maupun antarmuka pengguna dapat mengakses struktur papan secara terstruktur.

Atribut-atribut kelas \texttt{Board} dirangkum pada Tabel~\ref{tab:attr-board}.

\begin{longtable}{llp{6.5cm}}
	\caption{Atribut kelas \texttt{Board}}
	\label{tab:attr-board} \\
	\toprule
	\textbf{Atribut} & \textbf{Tipe} & \textbf{Deskripsi} \\
	\midrule
	\endfirsthead
	\toprule
	\textbf{Atribut} & \textbf{Tipe} & \textbf{Deskripsi} \\
	\midrule
	\endhead
	\bottomrule
	\endfoot
	\texttt{size} & \texttt{int} & Ukuran papan ($n$), menyatakan jumlah baris, kolom, dan sektor warna \\
	\midrule
	\texttt{colorMap} & \texttt{LinkedHashMap<String, List<int[]>{}>} & Pemetaan dari kunci warna RGB (format \texttt{"R,G,B"}) ke daftar koordinat $[\textit{baris}, \textit{kolom}]$ seluruh sel yang memiliki warna tersebut. Penggunaan \texttt{LinkedHashMap} mempertahankan urutan insersi warna, yang berpengaruh pada urutan pemrosesan variabel oleh \textit{solver} \\
	\midrule
	\texttt{colorSymbolMap} & \texttt{LinkedHashMap<String, String>} & Pemetaan dari kunci warna RGB ke simbol huruf (A, B, C, dst.) untuk keperluan visualisasi tekstual papan \\
	\bottomrule
\end{longtable}

Pemilihan \texttt{LinkedHashMap} sebagai struktur data untuk \texttt{colorMap} dan \texttt{colorSymbolMap}, dan bukan \texttt{HashMap} biasa, didasarkan pada kebutuhan akan urutan iterasi yang deterministik. Kedua kelas tersebut memiliki kompleksitas akses (\textit{lookup}) yang setara, yaitu $O(1)$ rata-rata. Perbedaannya terletak pada urutan iterasi: \texttt{HashMap} mengiterasi entri-entri berdasarkan urutan internal \textit{hash bucket} yang tidak terdefinisi dan dapat berubah antarsesi eksekusi, sedangkan \texttt{LinkedHashMap} mempertahankan urutan insersi entri. Pada konteks \textit{solver}, urutan iterasi warna pada \texttt{colorMap} secara langsung menentukan urutan pemrosesan variabel oleh algoritma \textit{backtracking}. Apabila urutan tersebut berubah antarsesi, perilaku \textit{solver} (jumlah langkah, jumlah \textit{backtrack}, dan cabang pencarian yang dieksplorasi) juga akan berubah, sehingga hasil eksperimen menjadi tidak \textit{reproducible}. Dengan \texttt{LinkedHashMap}, urutan tersebut konsisten antarsesi, sehingga eksperimen dapat direplikasi secara persis dan perilaku \textit{solver} dapat dianalisis dengan andal.

Metode-metode kelas \texttt{Board} dirangkum pada Tabel~\ref{tab:method-board}.

\begin{longtable}{lp{8.5cm}}
	\caption{Metode kelas \texttt{Board}}
	\label{tab:method-board} \\
	\toprule
	\textbf{Metode} & \textbf{Deskripsi} \\
	\midrule
	\endfirsthead
	\toprule
	\textbf{Metode} & \textbf{Deskripsi} \\
	\midrule
	\endhead
	\bottomrule
	\endfoot
	\texttt{Board(int, List<Cell>)} & \textit{Constructor} yang menerima ukuran papan dan daftar objek \texttt{Cell}, kemudian memanggil \texttt{buildColorMap()} dan \texttt{buildColorSymbols()} untuk menginisialisasi kedua peta internal \\
	\midrule
	\texttt{buildColorMap(List<Cell>)} & Mengiterasi seluruh objek \texttt{Cell} dan mengelompokkan koordinat sel ke dalam \texttt{colorMap} berdasarkan kunci warna RGB \\
	\midrule
	\texttt{buildColorSymbols()} & Memetakan setiap kunci warna RGB yang unik ke sebuah simbol huruf berurutan (A, B, C, dst.) untuk representasi tekstual \\
	\midrule
	\texttt{getSize()} & Mengembalikan ukuran papan \\
	\midrule
	\texttt{getColorMap()} & Mengembalikan peta warna-ke-posisi \\
	\midrule
	\texttt{getCellsByColor(String)} & Mengembalikan daftar koordinat seluruh sel yang memiliki warna tertentu \\
	\midrule
	\texttt{getSymbolForColor(String)} & Mengembalikan simbol huruf yang merepresentasikan suatu kunci warna RGB \\
	\midrule
	\texttt{printSymbolBoard()} & Mencetak papan ke konsol dengan setiap sel ditampilkan sebagai simbol huruf warnanya \\
	\midrule
	\texttt{printColorSummary()} & Mencetak ringkasan distribusi warna papan, termasuk jumlah sel per warna dan total warna unik \\
	\bottomrule
\end{longtable}

Metode \texttt{buildColorMap()} merupakan metode inti yang dipanggil pada saat konstruksi objek. Metode ini mengiterasi seluruh objek \texttt{Cell} yang diterima dari \texttt{BoardImporter}, membentuk kunci warna berformat \texttt{"R,G,B"} dari komponen RGB setiap sel, dan memasukkan koordinat $[\textit{baris}, \textit{kolom}]$ sel tersebut ke dalam daftar yang bersesuaian pada \texttt{colorMap}. Hasil dari proses ini adalah sebuah struktur pemetaan di mana setiap sektor warna terasosiasi dengan daftar lengkap posisi sel-selnya, yang kemudian digunakan oleh kedua \textit{solver} sebagai definisi domain variabel CSP.

\subsubsection{Kelas \texttt{BoardImporter}}
\label{subsubsec:kelas-boardimporter}

Kelas \texttt{BoardImporter} bertanggung jawab membaca dan mengurai berkas JSON papan permainan menjadi daftar objek \texttt{Cell}. Kelas ini berfungsi sebagai jembatan antara data eksternal (berkas JSON) dengan representasi internal sistem (objek \texttt{Cell} dan \texttt{Board}).

Kelas ini tidak memiliki atribut \textit{instance} karena seluruh metodenya bersifat statis. Konstanta \texttt{BOARD\_FOLDER} menyimpan jalur direktori berkas papan. Metode-metode kelas \texttt{BoardImporter} dirangkum pada Tabel~\ref{tab:method-boardimporter}.

\begin{longtable}{lp{8.5cm}}
	\caption{Metode kelas \texttt{BoardImporter}}
	\label{tab:method-boardimporter} \\
	\toprule
	\textbf{Metode} & \textbf{Deskripsi} \\
	\midrule
	\endfirsthead
	\toprule
	\textbf{Metode} & \textbf{Deskripsi} \\
	\midrule
	\endhead
	\bottomrule
	\endfoot
	\texttt{importBoard(int, int)} & Menerima ukuran papan dan nomor level, membaca berkas JSON yang bersesuaian, mengurai setiap entri menjadi objek \texttt{Cell}, dan mengembalikan daftar seluruh sel \\
	\midrule
	\texttt{parseRgba(String)} & Mengurai \textit{string} warna dalam format \texttt{rgba(R,G,B,A)} menjadi array komponen RGB bertipe integer \\
	\midrule
	\texttt{printBoard(List<Cell>, int)} & Mencetak papan ke konsol dengan menampilkan nilai RGB setiap sel untuk keperluan \textit{debugging} \\
	\bottomrule
\end{longtable}

Metode \texttt{importBoard()} merupakan metode utama kelas ini. Metode ini mengonstruksi nama berkas berdasarkan format \texttt{\{ukuran\}x\{ukuran\}\_level\{nomor\}.json}, membaca isi berkas sebagai \textit{string}, kemudian menguraikan \textit{string} tersebut menjadi \texttt{JSONArray}. Untuk setiap objek JSON dalam \textit{array}, metode mengekstrak atribut \texttt{row}, \texttt{col}, dan \texttt{color}, memanggil \texttt{parseRgba()} untuk mengonversi \textit{string} warna menjadi komponen RGB, dan membentuk objek \texttt{Cell} yang ditambahkan ke daftar hasil.

\subsubsection{Kelas \texttt{BacktrackingSolverAC3}}
\label{subsubsec:kelas-backtracking}

Kelas \texttt{BacktrackingSolverAC3} mengimplementasikan \textit{solver} berbasis \textit{backtracking} yang diperkuat dengan \textit{forward checking} dan propagasi AC-3 untuk menyelesaikan permasalahan \textit{Colored Queens}. Kelas ini juga mendefinisikan dua kelas bantu internal: \texttt{PruneAction} yang mencatat informasi pemangkasan domain untuk keperluan pemulihan, dan \texttt{Pair<K,V>} yang digunakan sebagai representasi busur pada antrean AC-3.

Atribut-atribut kelas \texttt{BacktrackingSolverAC3} dirangkum pada Tabel~\ref{tab:attr-backtracking}.

\begin{longtable}{llp{5.5cm}}
	\caption{Atribut kelas \texttt{BacktrackingSolverAC3}}
	\label{tab:attr-backtracking} \\
	\toprule
	\textbf{Atribut} & \textbf{Tipe} & \textbf{Deskripsi} \\
	\midrule
	\endfirsthead
	\toprule
	\textbf{Atribut} & \textbf{Tipe} & \textbf{Deskripsi} \\
	\midrule
	\endhead
	\bottomrule
	\endfoot
	\texttt{board} & \texttt{Board} & Objek papan yang menjadi masukan \textit{solver} \\
	\midrule
	\texttt{size} & \texttt{int} & Ukuran papan ($n \times n$) \\
	\midrule
	\texttt{colorCells} & \texttt{Map<String, List<int[]>{}>} & Pemetaan warna ke daftar koordinat sel, diperoleh dari \texttt{Board} \\
	\midrule
	\texttt{colors} & \texttt{List<String>} & Daftar kunci warna yang menentukan urutan pemrosesan variabel \\
	\midrule
	\texttt{solution} & \texttt{int[]} & Array solusi di mana \texttt{solution[$i$]} menyimpan indeks sel terpilih dalam domain warna ke-$i$; bernilai $-1$ apabila belum ditugaskan \\
	\midrule
	\texttt{occupied} & \texttt{boolean[][]} & Matriks yang menandai posisi sel yang telah ditempati menteri \\
	\midrule
	\texttt{validCells} & \texttt{Map<String, BitSet>} & Peta dari kunci warna ke \texttt{BitSet} yang merepresentasikan status kevalidan setiap sel dalam domain; bit bernilai \texttt{true} apabila sel masih valid \\
	\midrule
	\texttt{colorCellCount} & \texttt{int[]} & Pencacah ukuran domain efektif setiap variabel warna \\
	\midrule
	\texttt{pruneStack} & \texttt{Stack<PruneAction>} & Tumpukan yang mencatat seluruh pemangkasan domain secara kronologis untuk mendukung pemulihan saat \textit{backtrack} \\
	\midrule
	\texttt{steps} & \texttt{long} & Pencacah jumlah penugasan yang berhasil dilanjutkan ke tingkat rekursi berikutnya \\
	\midrule
	\texttt{backtracks} & \texttt{long} & Pencacah jumlah pembatalan penugasan \\
	\midrule
	\texttt{startTime} & \texttt{long} & Waktu mulai eksekusi dalam milidetik \\
	\bottomrule
\end{longtable}

Metode-metode kelas \texttt{BacktrackingSolverAC3} dirangkum pada Tabel~\ref{tab:method-backtracking}.

\begin{longtable}{lp{8.5cm}}
	\caption{Metode kelas \texttt{BacktrackingSolverAC3}}
	\label{tab:method-backtracking} \\
	\toprule
	\textbf{Metode} & \textbf{Deskripsi} \\
	\midrule
	\endfirsthead
	\toprule
	\textbf{Metode} & \textbf{Deskripsi} \\
	\midrule
	\endhead
	\bottomrule
	\endfoot
	\texttt{solve()} & Metode publik yang memulai proses pencarian dengan memanggil \texttt{placeQueens(0)} dan mencatat waktu eksekusi \\
	\midrule
	\texttt{placeQueens(int)} & Prosedur rekursif utama yang menugaskan menteri untuk setiap variabel warna secara berurutan \\
	\midrule
	\texttt{forwardCheck(int, int, int)} & Memangkas sel-sel yang berkonflik dengan menteri yang baru ditempatkan dari domain seluruh variabel yang belum ditugaskan \\
	\midrule
	\texttt{revise(int, int)} & Memeriksa apakah setiap sel dalam domain variabel tujuan memiliki \textit{support} yang konsisten di variabel sumber \\
	\midrule
	\texttt{propagateArcs(Queue)} & Menjalankan propagasi busur AC-3 untuk menegakkan konsistensi busur antara seluruh pasang variabel yang belum ditugaskan \\
	\midrule
	\texttt{conflicts(int[], int[])} & Mengembalikan \texttt{true} apabila dua sel berkonflik (baris sama, kolom sama, atau bertetangga) \\
	\midrule
	\texttt{undoPrunes(int)} & Memulihkan seluruh pemangkasan yang dicatat pada \texttt{pruneStack} sejak posisi \textit{checkpoint} yang diberikan \\
	\midrule
	\texttt{anyColorExhausted(int)} & Memeriksa apakah terdapat variabel yang belum ditugaskan dengan domain berukuran nol \\
	\midrule
	\texttt{getSolutionAsGrid()} & Mengonversi solusi berbasis indeks menjadi matriks koordinat $[\textit{baris}, \textit{kolom}]$ untuk digunakan oleh antarmuka \\
	\midrule
	\texttt{getSteps()}, \texttt{getBacktracks()} & Mengembalikan metrik pencarian \\
	\bottomrule
\end{longtable}

Representasi domain setiap variabel pada \textit{solver backtracking} menggunakan \texttt{BitSet}, yaitu struktur data Java yang menyimpan sekumpulan nilai biner (\texttt{true}/\texttt{false}) secara padat dalam bentuk \textit{array of bits}. Setiap warna memiliki satu \texttt{BitSet} dengan panjang yang sama dengan jumlah sel pada domain awalnya, di mana bit pada indeks ke-$i$ bernilai \texttt{true} apabila sel ke-$i$ masih valid sebagai kandidat penempatan, dan \texttt{false} apabila sel tersebut telah dipangkas oleh \textit{forward checking} atau AC-3. Pemilihan \texttt{BitSet} dibandingkan struktur alternatif seperti \texttt{List<Integer>} atau \texttt{boolean[]} didasarkan pada tiga keunggulan utama. Pertama, operasi pemangkasan dan pemulihan domain dapat dilakukan dalam waktu $O(1)$ melalui metode \texttt{clear(i)} dan \texttt{set(i)}; pendekatan berbasis \texttt{List} memerlukan penghapusan dari posisi tengah yang berbiaya $O(n)$ dan menyulitkan pemulihan saat \textit{backtrack}. Kedua, \texttt{BitSet} menyediakan metode \texttt{nextSetBit(n)} yang secara efisien mencari indeks bit \texttt{true} berikutnya, sehingga iterasi atas sel-sel yang masih valid dapat dilakukan tanpa memeriksa sel-sel yang telah dipangkas. Pendekatan berbasis \texttt{boolean[]} memerlukan perulangan eksplisit yang memeriksa setiap elemen. Ketiga, \texttt{BitSet} menggunakan memori secara efisien dengan menyimpan 64 bit per \textit{long}, dibandingkan dengan satu \textit{byte} per elemen pada \texttt{boolean[]}. Mekanisme pencatatan dan pemulihan pemangkasan melalui \texttt{pruneStack} bekerja secara sinergis dengan \texttt{BitSet}: setiap kali sebuah bit di-\texttt{clear} (pemangkasan), aksi tersebut dicatat pada \textit{stack} sebagai objek \texttt{PruneAction}, dan saat \textit{backtrack} dilakukan, aksi-aksi tersebut di-\textit{undo} dengan memanggil \texttt{set} pada bit-bit yang dipangkas.

Berikutnya, Metode \texttt{placeQueens()} merupakan inti dari kelas ini. Metode ini bekerja secara rekursif \textit{depth-first}, memproses variabel warna satu per satu sesuai urutan yang ditentukan oleh \texttt{colors}. Untuk setiap sel valid dalam domain variabel saat ini, metode ini menempatkan menteri, menjalankan \textit{forward checking} dan propagasi AC-3, memeriksa apakah terdapat domain yang kosong, dan apabila tidak ada, melanjutkan pencarian ke variabel berikutnya secara rekursif. Apabila pencarian gagal, seluruh pemangkasan yang telah dilakukan dipulihkan melalui mekanisme \texttt{pruneStack} dan sel berikutnya dicoba. Pseudocode metode ini ditunjukkan pada Algoritma~\ref{alg:backtracking}.

\begin{algorithm}[!htbp]
	\caption{Prosedur rekursif \texttt{placeQueens} pada kelas \texttt{BacktrackingSolverAC3}}
	\label{alg:backtracking}
	\begin{algorithmic}[1]
		\Function{PlaceQueens}{$colorIndex$}
		\If{$colorIndex = |C|$}
		\State \Return \textbf{true} \Comment{Seluruh variabel telah ditugaskan}
		\EndIf
		\State $color \gets C[colorIndex]$
		\If{$|D_{color}| = 0$}
		\State \Return \textbf{false} \Comment{Domain kosong, gagal}
		\EndIf
		\For{setiap sel $s$ dalam $D_{color}$ yang masih valid}
		\State Tempatkan menteri pada sel $s$
		\State $checkpoint \gets$ posisi \textit{pruneStack} saat ini
		\State \Call{ForwardCheck}{$s.row, s.col, colorIndex$}
		\State \Call{PropagateArcs}{antrean busur variabel yang belum ditugaskan}
		\If{tidak ada domain variabel lain yang kosong}
		\If{\Call{PlaceQueens}{$colorIndex + 1$}}
		\State \Return \textbf{true}
		\EndIf
		\EndIf
		\State \Call{UndoPrunes}{$checkpoint$}
		\State Batalkan penempatan menteri pada sel $s$
		\EndFor
		\State \Return \textbf{false}
		\EndFunction
	\end{algorithmic}
\end{algorithm}

Metode \texttt{forwardCheck()} dipanggil setiap kali sebuah menteri ditempatkan. Metode ini mengiterasi seluruh variabel warna yang belum ditugaskan dan memeriksa setiap sel valid dalam domain variabel tersebut terhadap posisi menteri yang baru ditempatkan. Sel yang berada pada baris atau kolom yang sama, atau yang bertetangga langsung dalam delapan arah, dihapus dari domain dengan menonaktifkan bit yang bersesuaian pada \texttt{BitSet}. Setiap pemangkasan dicatat pada \texttt{pruneStack}. Pseudocode metode ini tertera pada Algoritma~\ref{alg:fc-impl}.


\begin{algorithm}[!htbp]
	\caption{Prosedur \texttt{forwardCheck} pada kelas \texttt{BacktrackingSolverAC3}}
	\label{alg:fc-impl}
	\begin{algorithmic}[1]
		\Function{ForwardCheck}{$row, col, placedColorIdx$}
		\For{$colorIdx \gets placedColorIdx + 1$ \textbf{hingga} $|C| - 1$}
		\State $valid \gets$ \texttt{BitSet} domain warna $colorIdx$
		\For{setiap indeks $i$ yang masih aktif dalam $valid$}
		\State $(r_2, k_2) \gets$ koordinat sel ke-$i$ dari warna $colorIdx$
		\If{$row = r_2 \;\lor\; col = k_2 \;\lor\; (|row - r_2| \leq 1 \;\land\; |col - k_2| \leq 1)$}
		\State Hapus bit $i$ dari $valid$
		\State Kurangi $colorCellCount[colorIdx]$
		\State Catat $(colorIdx, i)$ pada \textit{pruneStack}
		\EndIf
		\EndFor
		\EndFor
		\EndFunction
	\end{algorithmic}
\end{algorithm}

Metode \texttt{revise()} merupakan komponen inti dari propagasi AC-3. Untuk setiap sel yang masih valid dalam domain variabel tujuan, metode ini memeriksa apakah terdapat setidaknya satu sel dalam domain variabel sumber yang tidak berkonflik (\textit{support}). Selain itu, sel tujuan juga diperiksa terhadap seluruh menteri yang telah ditempatkan. Apabila tidak ditemukan \textit{support} yang valid, sel tersebut dihapus dari domain dan dicatat pada \texttt{pruneStack}. Pseudocode metode ini diuraikan pada Algoritma~\ref{alg:revise-impl}.

\begin{algorithm}[!htbp]
	\caption{Fungsi \texttt{revise} pada kelas \texttt{BacktrackingSolverAC3}}
	\label{alg:revise-impl}
	\begin{algorithmic}[1]
		\Function{Revise}{$fromColorIdx, toColorIdx$}
		\State $revised \gets$ \textbf{false}
		\For{setiap indeks $i$ yang masih aktif dalam $D_{to}$}
		\State $supported \gets$ \textbf{false}
		\For{setiap indeks $j$ yang masih aktif dalam $D_{from}$}
		\If{$\lnot\;\text{konflik}(sel_j^{from},\; sel_i^{to})$}
		\State $supported \gets$ \textbf{true}
		\State \textbf{break}
		\EndIf
		\EndFor
		\If{$supported$}
		\For{setiap variabel $p$ yang telah ditugaskan, $p \neq toColorIdx$}
		\If{$\text{konflik}(sel_p^{placed},\; sel_i^{to})$}
		\State $supported \gets$ \textbf{false}
		\State \textbf{break}
		\EndIf
		\EndFor
		\EndIf
		\If{$\lnot\; supported$}
		\State Hapus bit $i$ dari $D_{to}$
		\State Kurangi $colorCellCount[toColorIdx]$
		\State Catat $(toColorIdx, i)$ pada \textit{pruneStack}
		\State $revised \gets$ \textbf{true}
		\EndIf
		\EndFor
		\State \Return $revised$
		\EndFunction
	\end{algorithmic}
\end{algorithm}

Metode \texttt{propagateArcs()} menjalankan propagasi busur AC-3 secara iteratif. Antrean diinisialisasi dengan seluruh pasang busur $(i, j)$ untuk $i \neq j$ di antara variabel yang belum ditugaskan. Untuk setiap busur yang diambil dari antrean, metode \texttt{revise()} dipanggil. Apabila domain variabel tujuan mengalami perubahan, busur-busur baru yang melibatkan variabel tersebut ditambahkan ke antrean untuk propagasi lebih lanjut. Propagasi berhenti apabila antrean kosong atau domain suatu variabel menjadi kosong. Pseudocode metode ini terdapat pada Algoritma~\ref{alg:propagate-impl}.

\begin{algorithm}[!htbp]
	\caption{Prosedur \texttt{propagateArcs} pada kelas \texttt{BacktrackingSolverAC3}}
	\label{alg:propagate-impl}
	\begin{algorithmic}[1]
		\Function{PropagateArcs}{$queue$}
		\While{$queue$ tidak kosong}
		\State $(from, to) \gets$ ambil busur dari $queue$
		\If{\Call{Revise}{$from, to$}}
		\If{$|D_{to}| = 0$}
		\State \Return \Comment{Domain kosong, hentikan propagasi}
		\EndIf
		\For{setiap variabel $k$, $k \neq to$, $k \neq from$}
		\State Tambahkan busur $(to, k)$ ke $queue$
		\EndFor
		\EndIf
		\EndWhile
		\EndFunction
	\end{algorithmic}
\end{algorithm}

\subsubsection{Kelas \texttt{Particle}}
\label{subsubsec:kelas-particle}

Kelas \texttt{Particle} merepresentasikan satu partikel dalam populasi PSO. Setiap partikel menyimpan sebuah kandidat solusi lengkap, vektor kecepatan yang diinterpretasikan sebagai probabilitas perubahan per dimensi, serta informasi \textit{personal best} yang dicapai selama proses optimasi.

Atribut-atribut kelas \texttt{Particle} dirangkum pada Tabel~\ref{tab:attr-particle}.

\begin{longtable}{llp{6cm}}
	\caption{Atribut kelas \texttt{Particle}}
	\label{tab:attr-particle} \\
	\toprule
	\textbf{Atribut} & \textbf{Tipe} & \textbf{Deskripsi} \\
	\midrule
	\endfirsthead
	\toprule
	\textbf{Atribut} & \textbf{Tipe} & \textbf{Deskripsi} \\
	\midrule
	\endhead
	\bottomrule
	\endfoot
	\texttt{candidate} & \texttt{int[]} & Posisi saat ini; \texttt{candidate[$i$]} menyimpan indeks sel terpilih dalam domain warna ke-$i$ \\
	\midrule
	\texttt{velocity} & \texttt{double[]} & Vektor kecepatan; \texttt{velocity[$j$]} $\in [0, 1]$ menyatakan probabilitas perubahan posisi pada dimensi $j$ \\
	\midrule
	\texttt{fitness} & \texttt{double} & Nilai \textit{fitness} saat ini, dihitung dari jumlah pelanggaran kendala berbobot \\
	\midrule
	\texttt{pBest} & \texttt{int[]} & Posisi terbaik yang pernah dicapai oleh partikel ini \\
	\midrule
	\texttt{pBestFitness} & \texttt{double} & Nilai \textit{fitness} dari posisi \texttt{pBest} \\
	\bottomrule
\end{longtable}

Kelas ini menyediakan \textit{constructor} yang menerima array posisi awal dan array kecepatan awal, serta metode \textit{getter} dan \textit{setter} untuk seluruh atribut. Kelas \texttt{Particle} tidak mengandung logika komputasi; seluruh perhitungan \textit{fitness} dan pembaruan posisi dikelola oleh kelas \texttt{PSOSolver}.

\subsubsection{Kelas \texttt{PSOSolver}}
\label{subsubsec:kelas-psosolver}

Kelas \texttt{PSOSolver} mengimplementasikan \textit{solver} berbasis \textit{Particle Swarm Optimization} diskrit. Kelas ini mengelola populasi partikel, topologi \textit{neighborhood}, serta seluruh mekanisme pembaruan kecepatan, posisi, dan \textit{fitness}.

Atribut-atribut kelas \texttt{PSOSolver} dirangkum pada Tabel~\ref{tab:attr-pso}.

\begin{longtable}{llp{5.5cm}}
	\caption{Atribut kelas \texttt{PSOSolver}}
	\label{tab:attr-pso} \\
	\toprule
	\textbf{Atribut} & \textbf{Tipe} & \textbf{Deskripsi} \\
	\midrule
	\endfirsthead
	\toprule
	\textbf{Atribut} & \textbf{Tipe} & \textbf{Deskripsi} \\
	\midrule
	\endhead
	\bottomrule
	\endfoot
	\texttt{board} & \texttt{Board} & Objek papan yang menjadi masukan \textit{solver} \\
	\midrule
	\texttt{size} & \texttt{int} & Ukuran papan \\
	\midrule
	\texttt{colorCells} & \texttt{Map<String, List<int[]>{}>} & Pemetaan warna ke daftar koordinat sel \\
	\midrule
	\texttt{colors} & \texttt{List<String>} & Daftar kunci warna \\
	\midrule
	\texttt{solution} & \texttt{int[]} & Array solusi \\
	\midrule
	\texttt{particles} & \texttt{List<Particle>} & Populasi partikel \\
	\midrule
	\texttt{neighborhoods} & \texttt{List<List<Integer>{}>} & Partisi indeks partikel ke dalam kelompok \textit{neighborhood} \\
	\midrule
	\texttt{nBests} & \texttt{int[]} & Indeks partikel terbaik dalam setiap \textit{neighborhood} \\
	\midrule
	\texttt{nIterations} & \texttt{int} & Batas maksimum iterasi ($I_{\max}$) \\
	\midrule
	\texttt{nParticles} & \texttt{int} & Jumlah partikel ($P$) \\
	\midrule
	\texttt{nNeighborhood} & \texttt{int} & Jumlah kelompok \textit{neighborhood} ($N_{\textit{nh}}$) \\
	\midrule
	\texttt{c1}, \texttt{c2} & \texttt{double} & Koefisien akselerasi kognitif dan sosial \\
	\midrule
	\texttt{inertia} & \texttt{double} & Koefisien inersia ($w$) \\
	\midrule
	\texttt{w1}, \texttt{w2} & \texttt{double} & Bobot komponen \textit{attacking} dan \textit{adjacency} pada fungsi \textit{fitness} \\
	\midrule
	\texttt{maxStagnation} & \texttt{int} & Batas iterasi tanpa perbaikan sebelum terminasi dini ($S_{\max}$) \\
	\midrule
	\texttt{random} & \texttt{Random} & Generator bilangan acak \\
	\bottomrule
\end{longtable}

Metode-metode kelas \texttt{PSOSolver} dirangkum pada Tabel~\ref{tab:method-pso}.

\begin{longtable}{>{\raggedright\arraybackslash}p{6cm} >{\raggedright\arraybackslash}p{7.5cm}}
	\caption{Metode kelas \texttt{PSOSolver}}
	\label{tab:method-pso} \\
	\toprule
	\textbf{Metode} & \textbf{Deskripsi} \\
	\midrule
	\endfirsthead
	\toprule
	\textbf{Metode} & \textbf{Deskripsi} \\
	\midrule
	\endhead
	\bottomrule
	\endfoot
	\texttt{solve()} & Metode publik utama yang menjalankan keseluruhan proses optimasi PSO \\
	\midrule
	\texttt{initialize()} & Menjalankan tiga tahap inisialisasi: pembangkitan partikel, pembentukan \textit{neighborhood}, dan evaluasi \textit{fitness} awal \\
	\midrule
	\texttt{generateParticles()} & Membangkitkan $P$ partikel dengan posisi dan kecepatan acak \\
	\midrule
	\texttt{generateNeighborhoods()} & Mempartisi seluruh partikel ke dalam $N_{\textit{nh}}$ kelompok \textit{neighborhood} dengan ukuran seimbang \\
	\midrule
	\texttt{calculateInitialFitness()} & Mengevaluasi \textit{fitness} awal seluruh partikel dan menentukan \texttt{nBest} setiap \textit{neighborhood} \\
	\midrule
	\texttt{calculateFitness(Particle)} & Menghitung nilai \textit{fitness} partikel berdasarkan pelanggaran baris/kolom dan \textit{adjacency} berbobot \\
	\midrule
	\texttt{calculateAttackingViolation}\newline\texttt{(Particle)} & Menghitung jumlah pasangan menteri yang berada pada baris atau kolom yang sama \\
	\midrule
	\texttt{calculateAdjacencyViolation}\newline\texttt{(Particle, boolean[][])} & Menghitung jumlah pasangan menteri yang bertetangga langsung dalam delapan arah \\
	\midrule
	\texttt{updateVelocity()} & Memperbarui kecepatan seluruh partikel berdasarkan formula PSO diskrit \\
	\midrule
	\texttt{updateParticles()} & Memperbarui posisi seluruh partikel berdasarkan kecepatan, mengevaluasi ulang \textit{fitness}, dan memperbarui \texttt{pBest} serta \texttt{nBest} \\
	\midrule
	\texttt{checkNbest()} & Memeriksa apakah salah satu \texttt{nBest} memiliki \textit{fitness} bernilai nol (solusi valid ditemukan) \\
	\midrule
	\texttt{checkLowestNBest()} & Mengembalikan indeks partikel dengan \textit{fitness} terendah di antara seluruh \texttt{nBest} \\
	\midrule
	\texttt{getLowestNBestFitness()} & Mengembalikan nilai \textit{fitness} terendah di antara seluruh \texttt{nBest} untuk deteksi stagnasi \\
	\bottomrule
\end{longtable}

Metode \texttt{solve()} merupakan metode utama yang mengorkestrasi keseluruhan proses PSO. Setelah inisialisasi populasi, metode ini menjalankan \textit{loop} utama yang pada setiap iterasi memperbarui kecepatan seluruh partikel, memperbarui posisi dan mengevaluasi ulang \textit{fitness}, serta memeriksa tiga kondisi terminasi: ditemukannya solusi valid, tercapainya batas iterasi, atau terdeteksinya stagnasi. Pseudocode metode ini dapat dilihat pada Algoritma~\ref{alg:pso-main}.

\begin{algorithm}[!htbp]
	\caption{Prosedur \texttt{solve} pada kelas \texttt{PSOSolver}}
	\label{alg:pso-main}
	\begin{algorithmic}[1]
		\Function{Solve}{}
		\State Inisialisasi populasi, \textit{neighborhood}, dan \textit{fitness} awal
		\State $f_{\textit{globalBest}} \gets$ \textit{fitness} terendah di antara seluruh \texttt{nBest}
		\State $stagnasi \gets 0$
		\For{$iter \gets 1$ \textbf{hingga} $I_{\max}$}
		\State Perbarui kecepatan seluruh partikel
		\State Perbarui posisi dan evaluasi \textit{fitness} seluruh partikel
		\If{terdapat \texttt{nBest} dengan \textit{fitness} $= 0$}
		\State \Return solusi valid ditemukan
		\EndIf
		\State $f_{\textit{current}} \gets$ \textit{fitness} terendah di antara seluruh \texttt{nBest}
		\If{$f_{\textit{current}} < f_{\textit{globalBest}}$}
		\State $f_{\textit{globalBest}} \gets f_{\textit{current}}$
		\State $stagnasi \gets 0$
		\Else
		\State $stagnasi \gets stagnasi + 1$
		\EndIf
		\If{$stagnasi \geq S_{\max}$}
		\State \Return terminasi dini, laporkan solusi terbaik
		\EndIf
		\EndFor
		\State \Return batas iterasi tercapai, laporkan solusi terbaik
		\EndFunction
	\end{algorithmic}
\end{algorithm}

Metode \texttt{updateVelocity()} dan \texttt{updateParticles()} bersama-sama mengimplementasikan mekanisme eksplorasi PSO diskrit. Pada \texttt{updateVelocity()}, kecepatan baru setiap dimensi dihitung menggunakan komponen inersia, kognitif, dan sosial dengan indikator perbedaan biner, kemudian di-\textit{clamp} pada interval $[0, 1]$. Pada \texttt{updateParticles()}, kecepatan digunakan sebagai probabilitas untuk mengadopsi posisi \textit{neighborhood best} pada setiap dimensi, kemudian \textit{fitness} dievaluasi ulang dan \texttt{pBest} serta \texttt{nBest} diperbarui apabila ditemukan nilai yang lebih baik. Pseudocode gabungan kedua metode ini disajikan pada Algoritma~\ref{alg:pso-update}.

\begin{algorithm}[!htbp]
	\caption{Prosedur \texttt{updateVelocity} dan \texttt{updateParticles} pada kelas \texttt{PSOSolver}}
	\label{alg:pso-update}
	\begin{algorithmic}[1]
		\For{setiap \textit{neighborhood} $N_i$}
		\State $\mathbf{x}_{\textit{nBest}} \gets$ posisi \texttt{nBest} dari $N_i$
		\For{setiap partikel $p$ dalam $N_i$}
		\State \Comment{\textit{Pembaruan kecepatan}}
		\For{$j \gets 0$ \textbf{hingga} $n - 1$}
		\State $r_1, r_2 \gets$ bilangan acak $\in [0, 1)$
		\State $\delta_{\textit{pBest}} \gets \mathbf{1}[x_j \neq \textit{pBest}_j]$
		\State $\delta_{\textit{nBest}} \gets \mathbf{1}[x_j \neq \textit{nBest}_j]$
		\State $v_j \gets \text{clamp}(w \cdot v_j + c_1 \cdot r_1 \cdot \delta_{\textit{pBest}} + c_2 \cdot r_2 \cdot \delta_{\textit{nBest}},\; 0,\; 1)$
		\EndFor
		\State \Comment{\textit{Pembaruan posisi}}
		\For{$j \gets 0$ \textbf{hingga} $n - 1$}
		\State $r \gets$ bilangan acak $\in [0, 1)$
		\If{$r < v_j$}
		\State $x_j \gets \textit{nBest}_j$
		\EndIf
		\EndFor
		\State Evaluasi ulang \textit{fitness} partikel $p$
		\State Perbarui \texttt{pBest} dan \texttt{nBest} jika diperlukan
		\EndFor
		\EndFor
	\end{algorithmic}
\end{algorithm}

\subsubsection{Kelas \texttt{Main}}
\label{subsubsec:kelas-main}

Kelas \texttt{Main} berfungsi sebagai titik masuk aplikasi. Kelas ini hanya berisi satu metode \texttt{main()} yang menginisialisasi objek \texttt{Homepage} pada \textit{Event Dispatch Thread} melalui \texttt{SwingUtilities.invokeLater()} untuk memastikan seluruh operasi antarmuka grafis berjalan pada \textit{thread} yang tepat sesuai konvensi Java Swing.

\subsubsection{Kelas \texttt{Homepage}}
\label{subsubsec:kelas-homepage}

Kelas \texttt{Homepage} merupakan subkelas dari \texttt{JFrame} yang merepresentasikan halaman utama aplikasi. Halaman ini menampilkan judul permainan, ikon menteri, dan tombol \textit{Start Game}. Kelas ini tidak memuat logika permainan; fungsi utamanya adalah menyediakan titik masuk visual bagi pengguna. Ketika tombol \textit{Start Game} ditekan, kelas ini membuat \textit{instance} baru dari \texttt{LevelSelectorWindow} dan menutup jendela saat ini.

\subsubsection{Kelas \texttt{LevelSelectorWindow}}
\label{subsubsec:kelas-levelselector}

Kelas \texttt{LevelSelectorWindow} merupakan subkelas dari \texttt{JFrame} yang menyajikan halaman pemilihan level. Halaman ini menampilkan pilihan ukuran papan yang dikelompokkan berdasarkan tingkat kesulitan (\textit{Easy}, \textit{Medium}, \textit{Hard}, dan \textit{Challenge}). Setiap tombol ukuran papan, ketika ditekan, membuat \textit{instance} baru dari \texttt{GameWindow} dengan ukuran dan level awal yang bersesuaian.

\subsubsection{Kelas \texttt{GameWindow}}
\label{subsubsec:kelas-gamewindow}

Kelas \texttt{GameWindow} merupakan subkelas dari \texttt{JFrame} yang merepresentasikan halaman permainan utama. Kelas ini mengintegrasikan seluruh komponen permainan: rendering papan, interaksi pengguna, eksekusi \textit{solver} secara asinkron, pemeriksaan pelanggaran, fitur \textit{hint}, dan fitur \textit{auto-solve}. Kelas ini juga mendefinisikan kelas dalam \texttt{BoardPanel} yang bertanggung jawab menggambar papan permainan.

Atribut-atribut utama kelas \texttt{GameWindow} dirangkum pada Tabel~\ref{tab:attr-gamewindow}.

\begin{longtable}{llp{5.5cm}}
	\caption{Atribut utama kelas \texttt{GameWindow}}
	\label{tab:attr-gamewindow} \\
	\toprule
	\textbf{Atribut} & \textbf{Tipe} & \textbf{Deskripsi} \\
	\midrule
	\endfirsthead
	\toprule
	\textbf{Atribut} & \textbf{Tipe} & \textbf{Deskripsi} \\
	\midrule
	\endhead
	\bottomrule
	\endfoot
	\texttt{currentSize} & \texttt{int} & Ukuran papan yang sedang dimainkan \\
	\midrule
	\texttt{currentLevel} & \texttt{int} & Nomor level yang sedang dimainkan \\
	\midrule
	\texttt{board} & \texttt{Board} & Objek papan yang sedang aktif \\
	\midrule
	\texttt{solution} & \texttt{int[][]} & Solusi yang dihitung oleh \textit{solver}, digunakan oleh fitur \textit{hint} dan \textit{auto-solve} \\
	\midrule
	\texttt{queenPlaced} & \texttt{boolean[][]} & Matriks yang menandai posisi sel tempat pengguna telah menempatkan menteri \\
	\midrule
	\texttt{currentWorker} & \texttt{SwingWorker} & \textit{Worker thread} yang menjalankan \textit{solver} secara asinkron di latar belakang \\
	\midrule
	\texttt{gameTimer} & \texttt{Timer} & Pencacah waktu permainan (\textit{stopwatch}) \\
	\midrule
	\texttt{boardPanel} & \texttt{BoardPanel} & Panel yang merender papan permainan \\
	\bottomrule
\end{longtable}

\texttt{SwingWorker} adalah kelas utilitas pada Java Swing yang dirancang untuk menjalankan tugas berdurasi panjang (\textit{long-running task}) pada \textit{thread} terpisah tanpa membekukan antarmuka pengguna. Pada arsitektur Swing, seluruh pembaruan antarmuka dilakukan pada satu \textit{thread} khusus yang disebut \textit{Event Dispatch Thread} (EDT). Apabila sebuah komputasi berat dijalankan langsung pada EDT, seluruh antarmuka akan menjadi tidak responsif selama komputasi tersebut berlangsung: tombol tidak dapat ditekan, jendela tidak dapat di-\textit{repaint}, dan input pengguna terabaikan. \texttt{SwingWorker} mengatasi permasalahan ini dengan menyediakan dua metode utama: \texttt{doInBackground()} yang dieksekusi pada \textit{thread} latar belakang dan berisi komputasi berat, serta \texttt{done()} yang dieksekusi kembali pada EDT setelah \texttt{doInBackground()} selesai dan digunakan untuk memperbarui antarmuka berdasarkan hasil komputasi. Pada aplikasi yang dibangun, solver \textit{backtracking} dijalankan di dalam \texttt{doInBackground()}, sehingga proses pencarian solusi tidak menghambat interaksi pengguna dengan papan permainan. Saat solusi ditemukan, \texttt{done()} dipanggil secara aman pada EDT untuk menyimpan solusi tersebut ke dalam variabel \texttt{solution} yang kemudian digunakan oleh fitur \textit{hint} dan \textit{auto-solve}.

Metode-metode utama kelas \texttt{GameWindow} mencakup \texttt{loadBoard()} yang membaca berkas papan dan menjalankan \textit{solver} secara asinkron, \texttt{checkViolations()} yang memeriksa seluruh pasangan menteri terhadap kendala dan menandai pelanggaran secara visual, serta \texttt{giveHint()} yang membandingkan penempatan pengguna dengan solusi dan memberikan petunjuk visual. Kelas dalam \texttt{BoardPanel} menggambar tiga lapisan secara berurutan pada metode \texttt{paintComponent()}: latar belakang berwarna setiap sel, garis grid pembatas, dan ikon menteri pada posisi yang telah ditempati.

\subsubsection{Kelas \texttt{UITheme}}
\label{subsubsec:kelas-uitheme}

Kelas \texttt{UITheme} menyediakan konstanta visual dan komponen antarmuka yang digunakan secara konsisten di seluruh halaman aplikasi. Kelas ini mendefinisikan palet warna (seperti \texttt{CREAM}, \texttt{PLUM}, \texttt{ROSE}, \texttt{TAUPE}, \texttt{TEXT\_DARK}, dan \texttt{TEXT\_MUTED}), metode utilitas untuk membuat objek \texttt{Font} dengan tipografi yang konsisten, metode \texttt{roundedButton()} untuk membuat tombol dengan sudut membulat, serta kelas dalam \texttt{GradientPanel} yang merender latar belakang gradien pada setiap halaman.

\subsection{Antarmuka Pengguna}
\label{subsec:antarmuka}

Antarmuka pengguna terdiri atas tiga layar utama yang dinavigasi secara berurutan. Alur interaksi pengguna dengan aplikasi diilustrasikan pada Gambar \ref{fig:alur-app}.

\begin{figure}[!htbp]
	\centering
	\includegraphics[width=0.85\linewidth]{flowchart_app.png}
	\caption{Alur interaksi pengguna pada aplikasi \textit{Colored Queens}.}
	\label{fig:alur-app}
\end{figure}

Berikut adalah penjelasan setiap proses pada Gambar \ref{fig:alur-app}:

\begin{enumerate}
	\item \textbf{\textit{Welcome Page}.} Layar pertama (\texttt{Homepage}) menampilkan judul permainan, ikon menteri, dan tombol \textit{Start Game} yang mengarahkan pengguna ke halaman pemilihan level. Halaman ini berfungsi sebagai titik masuk visual dan tidak memuat logika permainan.
	
	\item \textbf{\textit{Level Selector}.} Layar kedua (\texttt{LevelSelectorWindow}) menyajikan pilihan ukuran papan yang dikelompokkan berdasarkan tingkat kesulitan: \textit{Easy} ($7 \times 7$), \textit{Medium} ($8 \times 8$ dan $9 \times 9$), serta \textit{Hard} ($10 \times 10$, $11 \times 11$, dan $12 \times 12$). Selain itu, terdapat kategori \textit{Challenge} yang menyediakan papan berukuran $20 \times 20$ dan $30 \times 30$ dengan masing-masing satu level. Pemilihan ukuran papan langsung membuka halaman permainan pada level pertama dari ukuran tersebut.
	
	\item \textbf{Memuat papan dan menjalankan solver.} Setelah pengguna memilih level, kelas \texttt{GameWindow} membaca berkas JSON melalui \texttt{BoardImporter}, mengonstruksi objek \texttt{Board}, dan menjalankan solver \textit{backtracking} secara asinkron menggunakan \texttt{SwingWorker}. Solver berjalan di latar belakang agar antarmuka tetap responsif selama proses pencarian berlangsung seperti yang sudah dijelaskan pada Bagian \ref{subsubsec:kelas-gamewindow}. Untuk papan berukuran $20 \times 20$ dan $30 \times 30$, solusi disimpan secara \textit{hardcoded} karena waktu komputasi solver pada ukuran tersebut terlalu lama untuk dieksekusi secara langsung. Solusi yang dihasilkan solver disimpan dan digunakan oleh fitur \textit{hint} dan \textit{auto-solve}.
	
	\item \textbf{Menunggu aksi pemain.} Sistem memasuki kondisi menunggu, di mana papan ditampilkan sebagai grid berwarna di tengah jendela disertai panel kontrol di bagian bawah. Dari kondisi ini, pengguna dapat melakukan salah satu dari tiga aksi berikut.
	
	\item \textbf{Pemain menempatkan bidak menteri pada papan permainan.} Pengguna menempatkan atau menghapus menteri dengan mengklik sel pada papan. Pencacah waktu mulai berjalan saat pengguna menempatkan menteri pertama.
	
	\item \textbf{Pemeriksaan pelanggaran.} Setiap kali pengguna menempatkan atau menghapus menteri, sistem memeriksa seluruh pasangan menteri terhadap kendala baris, kolom, \textit{adjacency}, dan kesamaan warna. Menteri yang terlibat dalam pelanggaran ditandai dengan ikon berwarna merah.
	
	\item \textbf{Apakah semua kendala terpenuhi?} Sistem memeriksa apakah jumlah menteri yang ditempatkan sama dengan ukuran papan dan seluruh posisi sesuai dengan solusi. Apabila tidak, sistem kembali ke kondisi menunggu aksi pemain.
	
	\item \textbf{Pemain memilih opsi \textit{Hint}.} Sistem menampilkan petunjuk visual berupa sorotan hijau pada posisi yang benar dan sorotan oranye pada posisi yang salah untuk satu sektor warna. Setelah petunjuk ditampilkan, sistem kembali ke kondisi menunggu aksi pemain.
	
	\item \textbf{Pemain memilih opsi \textit{Solve}.} Solusi lengkap yang telah dihitung di latar belakang diterapkan secara otomatis pada papan.
	
	\item \textbf{\textit{Stopwatch} diberhentikan dan pesan kemenangan ditampilkan.} Pencacah waktu dihentikan dan dialog kemenangan ditampilkan kepada pengguna.
\end{enumerate}

Rancangan antarmuka untuk setiap layar beserta elemen-elemen interaktifnya dapat dilihat pada Gambar \ref{fig:mockup-homepage} hingga Gambar \ref{fig:mockup-win}.

% ===== HOMEPAGE =====
\begin{figure}[!htbp]
	\centering
	\includegraphics[width=0.75\linewidth]{mockup/homepage.png}
	\caption{Rancangan halaman utama (\textit{Homepage}).}
	\label{fig:mockup-homepage}
\end{figure}

Penjelasan elemen pada Gambar \ref{fig:mockup-homepage}:
\begin{enumerate}
	\item Judul aplikasi yang menampilkan nama permainan.
	\item Tombol \textit{Play} yang mengarahkan pengguna ke halaman pemilihan level.
\end{enumerate}

% ===== LEVEL SELECTOR =====
\begin{figure}[!htbp]
	\centering
	\includegraphics[width=0.75\linewidth]{mockup/levelSelector.png}
	\caption{Rancangan halaman pemilihan level (\textit{Level Selector}).}
	\label{fig:mockup-levelselector}
\end{figure}

Penjelasan elemen pada Gambar \ref{fig:mockup-levelselector}:
\begin{enumerate}
	\item Judul bagian pemilihan tingkat kesulitan.
	\item Judul bagian tantangan (\textit{Challenge}) untuk papan berukuran besar.
	\item Label kategori kesulitan (\textit{Easy}, \textit{Medium}, \textit{Hard}) yang mengelompokkan ukuran papan.
	\item Tombol ukuran papan yang dapat dipilih oleh pengguna untuk memulai permainan.
	\item Tombol ukuran papan pada kategori \textit{Challenge} ($20 \times 20$ dan $30 \times 30$).
\end{enumerate}

% ===== GAME WINDOW =====
\begin{figure}[!htbp]
	\centering
	\includegraphics[width=0.75\linewidth]{mockup/level.png}
	\caption{Rancangan halaman permainan (\textit{Game Window}).}
	\label{fig:mockup-gamewindow}
\end{figure}

Penjelasan elemen pada Gambar \ref{fig:mockup-gamewindow}:
\begin{enumerate}
	\item Label ukuran papan yang sedang dimainkan.
	\item Label nomor level yang sedang dimainkan.
	\item Papan permainan berupa grid berwarna tempat pengguna menempatkan bidak menteri.
	\item Bidak menteri yang telah ditempatkan pada papan. Ikon berwarna hitam digunakan pada sel terang dan ikon berwarna putih pada sel gelap, dengan pemilihan berdasarkan luminansi warna sel.
	\item \textit{Dropdown} ukuran papan untuk berpindah ke ukuran lain.
	\item \textit{Dropdown} level untuk berpindah ke level lain dalam ukuran yang sama.
	\item Pencacah waktu (\textit{stopwatch}) yang mulai berjalan saat pengguna menempatkan menteri pertama.
	\item Tombol \textit{Hint} yang menampilkan petunjuk visual.
	\item Tombol \textit{Solve} yang menerapkan solusi lengkap secara otomatis.
\end{enumerate}

% ===== SOLVED BY SOLVER =====
\begin{figure}[!htbp]
	\centering
	\includegraphics[width=0.75\linewidth]{mockup/solved_by_solver.png}
	\caption{Tampilan papan setelah solver menerapkan solusi lengkap.}
	\label{fig:mockup-solved}
\end{figure}

Penjelasan elemen pada Gambar \ref{fig:mockup-solved}:
\begin{enumerate}
	\item Seluruh bidak menteri ditempatkan pada posisi yang benar oleh solver. Seluruh menteri ditampilkan dengan warna hijau untuk menandakan solusi yang valid.
\end{enumerate}

% ===== WARNING: ADJACENT =====
\begin{figure}[!htbp]
	\centering
	\includegraphics[width=0.75\linewidth]{mockup/warning_adjacent.png}
	\caption{Tampilan peringatan pelanggaran kendala \textit{adjacency}.}
	\label{fig:mockup-warn-adj}
\end{figure}

Penjelasan elemen pada Gambar \ref{fig:mockup-warn-adj}:
\begin{enumerate}
	\item Dua bidak menteri yang bertetangga secara diagonal ditandai berwarna merah, menunjukkan pelanggaran kendala \textit{adjacency}. Sel-sel di sekitar kedua menteri yang berkonflik juga disorot untuk memperjelas area pelanggaran.
\end{enumerate}

% ===== WARNING: COLOR =====
\begin{figure}[!htbp]
	\centering
	\includegraphics[width=0.75\linewidth]{mockup/warning_color.png}
	\caption{Tampilan peringatan pelanggaran kendala kesamaan warna.}
	\label{fig:mockup-warn-color}
\end{figure}

Penjelasan elemen pada Gambar \ref{fig:mockup-warn-color}:
\begin{enumerate}
	\item Dua bidak menteri yang ditempatkan pada sektor warna yang sama ditandai berwarna merah, menunjukkan pelanggaran kendala \textit{one queen per color}.
\end{enumerate}

% ===== WARNING: HORIZONTAL =====
\begin{figure}[!htbp]
	\centering
	\includegraphics[width=0.75\linewidth]{mockup/warning_horizontal.png}
	\caption{Tampilan peringatan pelanggaran kendala baris.}
	\label{fig:mockup-warn-horiz}
\end{figure}

Penjelasan elemen pada Gambar \ref{fig:mockup-warn-horiz}:
\begin{enumerate}
	\item Dua bidak menteri yang berada pada baris yang sama ditandai berwarna merah. Seluruh sel pada baris tersebut disorot untuk memperjelas area pelanggaran.
\end{enumerate}

% ===== WARNING: VERTICAL =====
\begin{figure}[!htbp]
	\centering
	\includegraphics[width=0.75\linewidth]{mockup/warning_vertical.png}
	\caption{Tampilan peringatan pelanggaran kendala kolom.}
	\label{fig:mockup-warn-vert}
\end{figure}

Penjelasan elemen pada Gambar \ref{fig:mockup-warn-vert}:
\begin{enumerate}
	\item Dua bidak menteri yang berada pada kolom yang sama ditandai berwarna merah. Seluruh sel pada kolom tersebut disorot untuk memperjelas area pelanggaran.
\end{enumerate}

% ===== HINT 1: CORRECT POSITION =====
\begin{figure}[!htbp]
	\centering
	\includegraphics[width=0.75\linewidth]{mockup/hint1.png}
	\caption{Tampilan fitur \textit{hint}: sorotan posisi yang benar.}
	\label{fig:mockup-hint1}
\end{figure}

Penjelasan elemen pada Gambar \ref{fig:mockup-hint1}:
\begin{enumerate}
	\item Sorotan hijau pada sel yang merupakan posisi benar untuk sektor warna yang belum terisi, menunjukkan di mana menteri seharusnya ditempatkan.
\end{enumerate}

% ===== HINT 2: WRONG + CORRECT =====
\begin{figure}[!htbp]
	\centering
	\includegraphics[width=0.75\linewidth]{mockup/hint2.png}
	\caption{Tampilan fitur \textit{hint}: sorotan posisi salah dan posisi benar.}
	\label{fig:mockup-hint2}
\end{figure}

Penjelasan elemen pada Gambar \ref{fig:mockup-hint2}:
\begin{enumerate}
	\item Menteri yang ditempatkan pada posisi salah ditandai dengan sorotan merah, menunjukkan bahwa penempatan perlu dipindahkan.
	\item Sorotan oranye pada sel yang merupakan posisi benar untuk sektor warna tersebut, menunjukkan ke mana menteri seharusnya dipindahkan.
\end{enumerate}

% ===== WIN =====
\begin{figure}[!htbp]
	\centering
	\includegraphics[width=0.75\linewidth]{mockup/win.png}
	\caption{Rancangan tampilan kemenangan.}
	\label{fig:mockup-win}
\end{figure}

Penjelasan elemen pada Gambar \ref{fig:mockup-win}:
\begin{enumerate}
	\item Ikon trofi sebagai elemen visual perayaan kemenangan.
	\item Pesan ucapan selamat kepada pemain.
	\item Waktu penyelesaian yang dicatat oleh \textit{stopwatch} selama permainan berlangsung.
	\item Tombol \textit{Next} untuk melanjutkan ke level berikutnya.
	\item Tombol \textit{Level Picker} untuk kembali ke halaman pemilihan level.
	\item Latar belakang halaman permainan yang diredupkan (\textit{dimmed}) untuk memberikan fokus pada dialog kemenangan.
\end{enumerate}

\subsection{Teknologi yang Digunakan}
\label{subsec:teknologi}

Tabel \ref{tab:teknologi} merangkum teknologi dan pustaka yang digunakan dalam pengembangan sistem.

\begin{longtable}{lp{9cm}}
	\caption{Teknologi yang digunakan}
	\label{tab:teknologi} \\
	\toprule
	\textbf{Teknologi} & \textbf{Penggunaan} \\
	\midrule
	\endfirsthead
	\toprule
	\textbf{Teknologi} & \textbf{Penggunaan} \\
	\midrule
	\endhead
	\bottomrule
	\endfoot
	Java (JDK 23) & Bahasa pemrograman utama untuk seluruh modul \\
	\midrule
	Java Swing / JFrame & \textit{Framework} antarmuka grafis untuk rendering papan, navigasi, dan interaksi pengguna \\
	\midrule
	\texttt{org.json} & Pustaka untuk penguraian berkas JSON papan \\
	\midrule
	JSON & Format penyimpanan data papan (\texttt{NxN\_levelM.json}) \\
	\bottomrule
\end{longtable}

Pemilihan Java Swing sebagai \text it{framework} antarmuka didasarkan pada kemudahan implementasi dan ketersediaannya sebagai bagian dari \textit{Java Standard Edition} tanpa memerlukan dependensi eksternal tambahan. Meskipun Swing merupakan teknologi yang lebih lama dibandingkan alternatif seperti JavaFX, kemampuannya untuk merender grid berwarna dan menangani interaksi klik pengguna cukup memadai untuk kebutuhan visualisasi papan \textit{Colored Queens}.